\section{Discussion}\label{sec:discussion}

\subsection{Connections to prior work}\label{sec:connections}

Our framework unifies several threads from prior work.

\paragraph{Test-time compute scaling.}
Snell et al.~\cite{snell2024scaling} show that compute-optimal scaling
strategies are difficulty-dependent: easy problems saturate quickly, while hard
problems require substantially more compute.  Our results formalize this
observation by identifying \emph{what makes a problem hard} for verification:
proof complexity measures, particularly the composite measure~$C(F)$.
Theorem~\ref{thm:convergence} provides a formula relating the convergence step
to~$C(F)$, replacing the model-dependent pass@1 metric with a formula-intrinsic
quantity.

\paragraph{Verification versus solving difficulty.}
Shi and Jin~\cite{shi2025heimdall} demonstrate that verification difficulty
does not correlate with problem-solving difficulty.  Our framework explains this
distinction: solving difficulty tracks search complexity (how hard it is to find
a proof), while verification difficulty tracks proof complexity (how hard it is
to check a given proof).  These are related but distinct quantities---a problem
may have a short proof that is hard to find, or a long proof that is easy to
locate but hard to verify.

\paragraph{Width--size tradeoffs.}
Ben-Sasson and Wigderson's~\cite{bensasson2001short} width--size relation
$w(F \vdash \bot) \le w(F) + O(\sqrt{n \ln S(F)})$ parallels our finding that
width and length jointly predict convergence better than either alone.  The
composite measure~$C(F)$, which combines both, achieves $\rho = -0.991$, while
individual measures achieve at most $\rho = -0.840$.

\paragraph{Space and practical hardness.}
Jarvisa\"alo et al.~\cite{jarvisalo2012relating} show that space complexity
predicts SAT solver hardness better than length or width.  In our experiments,
space ranks second ($\rho = -0.782$) after length ($\rho = -0.840$) in
predicting convergence rate.  This suggests that LLM verification may weight
these measures differently than CDCL-based solvers, a distinction that merits
empirical investigation.

\paragraph{Information-theoretic analysis.}
Ton et al.~\cite{ton2025understanding} prove that once an LLM encounters an
unidentifiable sub-task in a CoT chain, all subsequent steps yield zero
information gain.  This ``information plateau'' corresponds to the convergence
ceiling~$p_{\max}(F)$ in our model: once verification accuracy saturates, no
additional CoT steps improve performance.  Moreover, their result strengthens
our error accumulation bound (Theorem~\ref{thm:error-accumulation}), since
correlated errors make the bound optimistic.

\paragraph{Typed chain-of-thought.}
Perrier~\cite{perrier2025typed} maps CoT reasoning to typed proof terms via the
Curry--Howard correspondence, providing a formal mechanism for measuring proof
structure in natural language reasoning.  The Minimal Path Size (MPS) metric
introduced there is related to our proof depth measure~$d(F)$.  Combining
Perrier's typed CoT framework with our complexity--convergence analysis could
yield a unified theory of CoT verification for natural mathematical proofs.

\subsection{Implications}\label{sec:implications}

Our results suggest several practical implications.

\paragraph{Compute allocation.}
Proof complexity measures can guide the allocation of test-time compute.  Given
a proof instance~$F$, one can compute (or estimate) $C(F)$ and use
Theorem~\ref{thm:convergence} to predict the number of CoT steps and
Corollary~\ref{cor:compute-tradeoff} to predict the number of self-consistency
chains needed.  This replaces the trial-and-error approach to test-time
compute allocation with a principled, a priori estimate.

\paragraph{Selective human review.}
Instances with $C(F) > C^*$ fall in the harder regime identified by the phase
transition (Proposition~\ref{prop:phase-transition}).  These instances may
benefit from human review or from structural decomposition into sub-proofs with
lower individual complexity, rather than brute-force application of more CoT
chains.

\paragraph{Fundamental limits.}
The exponential decay in accuracy with proof length
(Theorem~\ref{thm:error-accumulation}) implies that step-by-step verification
of proofs with super-polynomial length faces a fundamental barrier.  For
formula families like pigeonhole, where any resolution proof has exponential
length, CoT verification accuracy degrades to chance level for large instances
(Example~\ref{ex:php-accuracy}).  Overcoming this barrier likely requires
verification strategies that do not process each step independently, such as
hierarchical verification or proof compression.

\subsection{Limitations}\label{sec:limitations}

We note several limitations of this work.

First, our theoretical framework is validated through simulation rather than
direct LLM experiments.  While the simulation is calibrated to published
empirical data, the strong correlations obtained (particularly $\rho = -0.991$
for the composite measure) partly reflect the parametric model construction.
The qualitative predictions---the ordering of difficulty across families and
the existence of saturation behavior---are more robust than the specific
correlation magnitudes.

Second, the step-independence assumption in
Theorem~\ref{thm:error-accumulation} is a simplification.  Real CoT traces
exhibit sequential dependencies, and the information-theoretic analysis of Ton
et al.~\cite{ton2025understanding} suggests that errors propagate in a
correlated fashion.

Third, our complexity measures are defined for propositional resolution proofs.
Extending these to first-order logic or to natural language mathematical proofs
requires additional formalization, such as the typed CoT framework of
Perrier~\cite{perrier2025typed}.

Fourth, the convergence rate model~\eqref{eq:lambda-model} uses fixed weights
$w_i$ calibrated to aggregate statistics.  These weights may vary across model
architectures and scales.

\subsection{Open questions}\label{sec:open}

We conclude this section with several open questions.

\begin{conjecture}\label{conj:empirical}
For at least one family among Tseitin, pigeonhole, and pebbling formulas, the
Spearman rank correlation between $C(F)$ and the empirical convergence rate of
GPT-4 or Claude verification exceeds $-0.8$ in absolute value.
\end{conjecture}

\begin{conjecture}\label{conj:space}
For CoT verification by LLMs with bounded context windows, proof space
$\operatorname{Sp}(F)$ is a stronger predictor of verification difficulty than
proof length~$S(F)$.
\end{conjecture}

\begin{conjecture}\label{conj:sharp}
In the limit of large formula size~$n$ for random $k$-CNF at the
satisfiability threshold, the phase transition in verification accuracy at
$C^*$ becomes sharp (i.e., the accuracy function is discontinuous in the
limit).
\end{conjecture}
